\documentclass[10pt,a4paper]{amsart}
\usepackage[margin=19mm]{geometry}
\usepackage{amsmath,amssymb,mathtools}
\usepackage[scr=rsfs,cal=euler]{mathalfa}
\usepackage{xfrac}
\usepackage[unicode,hidelinks,bookmarks=false]{hyperref}
\newcommand{\ie}{i.e.\ }
\newcommand{\cf}{cf.\ }
\newcommand{\resp}{respectively\ }
\newcommand{\Subsection}{\subsection}
\providecommand{\nobreakdash}{\nobreak\hbox{-}}
\setlength{\emergencystretch}{2em}
\multlinegap0pt
\usepackage{bm}
\usepackage{bbm}
\usepackage{dsfont}
\usepackage{xspace}
\usepackage{cancel}
\usepackage{mathtools}
\usepackage{amsthm}
\makeatletter
\@ifundefined{theorem}{\newtheorem{theorem}{Theorem}}{}
\@ifundefined{lemma}{\newtheorem{lemma}[theorem]{Lemma}}{}
\makeatother
\newcommand{\NNN}{\mathbf{N}}
\newcommand{\QQ}{\mathbb{Q}}
\newcommand{\ZZ}{\mathbb{Z}}
\newcommand{\xx}{\mathbf{x}}
\newcommand{\yy}{\mathbf{y}}
\renewcommand{\geq}{\geqslant}
\renewcommand{\ge}{\geqslant}
\renewcommand{\leq}{\leqslant}
\renewcommand{\le}{\leqslant}
\renewcommand{\mod}{\mathrm{mod}\,}
\title[Liouville function, von Mangoldt function and norm forms at ]{Liouville function, von Mangoldt function and norm forms at random binary forms}
\author{Yijie Diao}
\date{Source: arXiv:2506.18065v2 (CC BY 4.0)}
\begin{document}
\maketitle

\noindent\emph{Source and licence.} Liouville function, von Mangoldt function and norm forms at random binary forms. Version arXiv:2506.18065v2 (CC BY 4.0), distributed under the Creative Commons Attribution 4.0 International licence, as recorded in the arXiv licence metadata for this exact version and shown on the abstract page https://arxiv.org/abs/2506.18065v2. Full rights notice: licenses/arxiv-2506.18065-CC-BY-4.0.txt.\par\smallskip

\noindent\textit{Editorial scope.} Counted unit: the Lemma whose source label is \texttt{None} in the article (source0.tex lines 1844--1849, source label \texttt{None}), delivered together with the article's own complete original proof of it (source0.tex lines 1850--1942, one \texttt{proof} environment of 93 lines). No mathematics is written, repaired or re-derived: statement and proof are the article's own text line for line. Required context retained uncounted: norm-f (lines 294--296); norm-g (lines 318--320); cond-WA-1 (lines 1840--1842). Every definition, display and label of those lines is retained unchanged. Omitted from this package and not redistributed: 1--293, 297--317, 321--1839, 1843--1843, 1943--2129. Nothing inside a retained excerpt is omitted or altered: every retained body is delivered exactly as the article's lines are written, so no omission receipt of any kind accompanies this package. Citations: the retained excerpts carry no citation command at all, so no bibliography entry of the article is transcribed and no reference is redistributed. Reference adaptation: none was needed and none was made; every reference inside the delivered unit resolves inside it, so no unresolved reference or citation is produced. Printed slips kept literal: no printed word, symbol, hypothesis or number of the counted statement or of its proof was altered. Layout and class: the article's own class is \texttt{amsart} with packages including amsfonts, amsthm, amsmath, amssymb, hyperref, bm, graphicx, xparse, geometry, mathrsfs, comment, graphics, aliascnt, enumerate; the wrapper is the lane's pinned \texttt{amsart} editorial wrapper at 10pt on A4, which restates the theorem-like environments the retained text needs and adds the article's own macro definitions that the retained text needs (\texttt{\textbackslash NNN}, \texttt{\textbackslash QQ}, \texttt{\textbackslash ZZ}, \texttt{\textbackslash ge}, \texttt{\textbackslash geq}, \texttt{\textbackslash le}, \texttt{\textbackslash leq}, \texttt{\textbackslash mod}, \texttt{\textbackslash xx}, \texttt{\textbackslash yy}), with extra wrapper packages (bm, bbm, dsfont, xspace, cancel, mathtools). The delivered numbering is the unit's own. Assets: no retained span contains a figure environment, a graphics-inclusion command, a PDF-page inclusion, a table or a TikZ picture; no image file is copied into this package, the package declares no asset dependency and no image file is redistributed with it. Licence: version arXiv:2506.18065v2 of this article is distributed under the Creative Commons Attribution 4.0 International licence, as recorded in the arXiv licence metadata for this exact version and shown on the abstract page https://arxiv.org/abs/2506.18065v2. A full rights notice is delivered at licenses/arxiv-2506.18065-CC-BY-4.0.txt. Characters: every retained excerpt is pure ASCII, so the delivered engine, XeLaTeX with its default Latin Modern text font, reports no missing character.
\begin{eqnarray}\label{norm-f}
	\mathbf{N}_K(\mathbf{y}) = f(t) \neq 0.
\end{eqnarray}

	\begin{eqnarray}\label{norm-g}
		\mathbf{N}_K(x_1, \dots, x_e) = \sum_{0 \leq i \leq d} c_i u^i v^{d-i},
	\end{eqnarray}

\begin{align}\label{cond-WA-1}
	|\yy - \yy^{(p)} |_p \leq p^{-\kappa}, \quad |t - t^{(p)}|_p \leq p^{-\kappa}.
\end{align}

\begin{lemma}
	The equation (\ref{norm-f}) satisfies weak approximation at $p$, if for any integer $k > 0$ and any $(\xx, u, v) \in (\ZZ / p^k \ZZ)^{e+2}$ satisfying 
	$$\NNN_K(\xx) \equiv g(u, v) \ (\mod \, p^k),$$ 
	there exists integer point $(\xx^{(\ZZ)}, u^{(\ZZ)}, v^{(\ZZ)}) \in \ZZ^{e+2}$ with $v^{(\ZZ)} \neq 0$ of (\ref{norm-g}), such that 
	$$(\xx^{(\ZZ)}, u^{(\ZZ)}, v^{(\ZZ)}) \equiv (\xx, u, v) \ (\mod \, p^k). $$
\end{lemma}

\begin{proof}
We fix a $\QQ_p$-point $(\yy^{(p)},t^{(p)})$ on~\eqref{norm-f} and an integer $\kappa\ge 1$. We write  
\[
t^{(p)}=\frac{u^{(p)}}{v^{(p)}},  \,\text{ where } \,
u^{(p)},v^{(p)}\in\ZZ_p,\;
v^{(p)}\neq 0,\;
\min\bigl(v_p(u^{(p)}),v_p(v^{(p)})\bigr)=0.
\]
We set
\begin{align}\label{vpp}
	r=v_p\!\bigl(v^{(p)}\bigr)\ge 0.
\end{align}
We also put  
\[
\xx^{(p)}=(v^{(p)})^{d/e}\,\yy^{(p)}\in\QQ_p^{\,e}.
\]
Let $s\ge 0$ be the smallest integer with  
\begin{align}\label{def-s}
	N^{(d/e)}\xx^{(p)}\in\ZZ_p^{\,e},\qquad \text{ where } N = p^s.
\end{align}
Now we define  
\begin{equation}\label{eq:k-choice}
k=\kappa + \frac{2(r+s)d}{e}.
\end{equation}
Let $(\xx, u, v) \in (\ZZ / p^k \ZZ)^{e+2}$ satisfying
\[
(\xx,u,v)\equiv\bigl(N^{d/e}\xx^{(p)},\,Nu^{(p)},\,Nv^{(p)}\bigr)\pmod{p^{\,k}}.
\]
By assumption, there exists $(\xx^{(\ZZ)},u^{(\ZZ)},v^{(\ZZ)})\in\ZZ^{e+2}$ solving~\eqref{norm-g} with $v^{(\ZZ)} \neq 0$ and  
\begin{equation}\label{eq:congruence}
(\xx^{(\ZZ)},u^{(\ZZ)},v^{(\ZZ)})\equiv\bigl(N^{d/e}\xx^{(p)},\,Nu^{(p)},\,Nv^{(p)}\bigr)\pmod{p^{\,k}}.
\end{equation}
We set  
\[
\yy=(v^{(\ZZ)})^{-d/e}\xx^{(\ZZ)},\qquad t=\frac{u^{(\ZZ)}}{v^{(\ZZ)}}.
\]
Then the rational point $(\yy,t)$ lies on~\eqref{norm-f}.

We have
$$t-t^{(p)}
 = \frac{u^{(\ZZ)}v^{(p)}-u^{(p)}v^{(\ZZ)}}{v^{(\ZZ)}v^{(p)}}.$$
From~\eqref{eq:congruence}, we know
\[
u^{(\ZZ)}v^{(p)}-u^{(p)}v^{(\ZZ)}
 =(u^{(\ZZ)}-Nu^{(p)})v^{(p)}-u^{(p)}\bigl(v^{(\ZZ)} - Nv^{(p)}\bigr)\equiv 0\pmod{p^{\,k}},
\]
so $v_p\bigl(u^{(\ZZ)}v^{(p)}-u^{(p)}v^{(\ZZ)}\bigr)\ge k$.
By (\ref{eq:k-choice}) we know 
$$k > r + s = v_p(Nv^{(p)}).$$ 
By (\ref{eq:congruence}) we know $|v^{(\ZZ)} - Nv^{(p)}|_p \le p^{-k}$, it follows that 
\begin{align}\label{vpz}
	v_p(v^{(\ZZ)}) = v_p(N v^{(p)}) =  r+s.
\end{align}
Therefore, we have  
\begin{align}\label{p-t}
	|t-t^{(p)}|_p
 =\left|\frac{u^{(\ZZ)}v^{(p)}-u^{(p)}v^{(\ZZ)}}{v^{(\ZZ)}v^{(p)}}\right|_p
 \le p^{-k+2r+s}\le p^{-\kappa}.
\end{align}

For $1\le i\le e$, we have
\begin{align*}
y_i-y_i^{(p)}
&=\frac{x_i^{(\ZZ)}}{(v^{(\ZZ)})^{d/e}}-\frac{x_i^{(p)}}{(v^{(p)})^{d/e}}\\[2pt]
&=\frac{\bigl(x_i^{(\ZZ)}-N^{d/e}x_i^{(p)}\bigr)(v^{(p)})^{d/e}
         -\bigl((v^{(\ZZ)})^{d/e}-(Nv^{(p)})^{d/e}\bigr)x_i^{(p)}}
       {(v^{(\ZZ)}v^{(p)})^{d/e}}.
\end{align*}
Using(\ref{vpp}) and \eqref{eq:congruence} we have  
\begin{align*}
v_p \bigl(x_i^{(\ZZ)}-N^{d/e}x_i^{(p)}\bigr)\ge k,& \quad
v_p(v^{(p)}) = r.
\end{align*}
By ~\eqref{eq:congruence} we know
$$
v_p \bigl((v^{(\ZZ)})^{d/e}-(Nv^{(p)})^{d/e}\bigr) \geq v_p \bigl(v^{(\ZZ)}-Nv^{(p)}\bigr) \geq   k.$$ 
By (\ref{def-s}), we have $v_p(x_i^{(p)}) \geq -sd/e.$
Hence the numerator has valuation at least $k-sd/e$, while by (\ref{vpp}) and (\ref{vpz}) we know 
\[
v_p\bigl((v^{(\ZZ)}v^{(p)})^{d/e}\bigr)=(2r+s)\,d/e.
\]
Therefore, it follows from \eqref{eq:k-choice} that
\[
v_p\bigl(y_i-y_i^{(p)}\bigr) \geq k - sd/e - (2r+s)d/e = k-(2r+2s)d/e\ge\kappa.
\]
Thus 
\begin{align}\label{p-y}
	|\yy-\yy^{(p)}|_p\le p^{-\kappa}
\end{align}

\medskip
Now (\ref{p-t}) and (\ref{p-y}) ensure that both inequalities in~\eqref{cond-WA-1} are satisfied. Hence weak approximation at $p$ holds.
\end{proof}

\end{document}
